\documentclass[10pt]{article}
\usepackage[T1]{fontenc}
\usepackage[utf8]{inputenc}
\usepackage{lmodern}
\usepackage[margin=0.62in]{geometry}
\usepackage{microtype}
\usepackage{amsmath,amssymb}
\usepackage{titlesec}
\usepackage[hidelinks]{hyperref}
\titlespacing*{\section}{0pt}{0.65em}{0.25em}

\begin{document}
\begin{center}
{\LARGE Literature Resolution of Remark 7.3\par}
\vspace{0.15em}
{\large Right Bregman Klee Sets without Full Domain\par}
\vspace{0.45em}
Status note for arXiv:0802.2322 \quad $\cdot$ \quad 9 August 2026
\end{center}
\vspace{0.25em}

\noindent\textbf{Status.} \texttt{literature\_already\_answered}.

\section*{Original question}
Bauschke, Wang, Ye, and Yuan assume that $f$ is a Legendre function on
$\mathbb R^J$, write $U=\operatorname{int}\operatorname{dom}f$, and take a
nonempty compact set $C\subset U$.  Their Theorem 7.2 proves that a right
Bregman Klee set is a singleton when $\operatorname{dom}f=\mathbb R^J$.
Remark 7.3 on PDF page 15 then asks whether the theorem remains true when the
full-domain assumption is dropped \cite{BWYY2008}.

Here ``right Bregman Klee'' means that, for every $x\in U$, the maximizer of
\[
  c\longmapsto D_f(x,c)
  \qquad(c\in C)
\]
is unique.

\section*{Explicit supporting answer}
Bauschke, Macklem, Sewell, and Wang explicitly identify ``[7, Remark 7.3]''
and state on PDF page 4 that they settle the question affirmatively without
full domain.  Their Theorem 3.2 on PDF page 10 says \cite{BMSW2009}:
\begin{quote}
Suppose that $C$ is compact and that $C$ is right Bregman Klee.  Then $C$ is
a singleton.
\end{quote}
This has exactly the hypotheses left after removing
$\operatorname{dom}f=\mathbb R^J$ from the source theorem, so it is a direct
and complete answer.

\section*{Proof mechanism}
The later paper considers the convex farthest-distance function
\[
  F_C(x)=\sup_{c\in C}D_f(x,c).
\]
Its Theorem 3.1 proves that $F_C$ has a minimizer $x_0$ in $U$, even when
$\operatorname{dom}f$ has a boundary.  The proof uses coercivity of each
$D_f(\cdot,c)$ and a Bregman projection argument to rule out a boundary
minimizer.  At $x_0$, the Ioffe--Tikhomirov subdifferential formula and the
unique active farthest point $Q_C(x_0)$ give
\[
  0\in\partial F_C(x_0)
  =\{\nabla f(x_0)-\nabla f(Q_C(x_0))\}.
\]
The Legendre gradient is injective, hence $x_0=Q_C(x_0)$.  Thus the largest
value of $D_f(x_0,c)$ over $C$ is $D_f(x_0,x_0)=0$; nonnegativity and strict
convexity force $C=\{x_0\}$.

\section*{Scope}
The answer covers general finite-dimensional Legendre functions in the later
paper's standing setting, including non-full-domain examples such as
Kullback--Leibler and Itakura--Saito divergences.  It does not settle the
classical infinite-dimensional Hilbert-space farthest-point conjecture.

\begin{thebibliography}{9}
\bibitem{BWYY2008}
H. H. Bauschke, X. Wang, J. Ye, and X. Yuan,
\emph{Bregman distances and Klee sets},
\href{https://arxiv.org/abs/0802.2322}{arXiv:0802.2322}.

\bibitem{BMSW2009}
H. H. Bauschke, M. S. Macklem, J. B. Sewell, and X. Wang,
\emph{Klee sets and Chebyshev centers for the right Bregman distance},
\href{https://arxiv.org/abs/0908.2013}{arXiv:0908.2013}.
\end{thebibliography}

\end{document}
